\documentclass[a4paper, 11px]{article}

\usepackage[french]{babel}
\usepackage[utf8]{inputenc}
\usepackage{fancyhdr}
\usepackage{lastpage}
\usepackage{graphicx}
\usepackage{rotating}
\usepackage{textcomp}
\usepackage{xspace}
\usepackage[toc,page]{appendix}
\usepackage{array}
\usepackage{amssymb}
\usepackage{enumerate}
\usepackage{enumitem}
\usepackage{eso-pic}
\usepackage{hyperref}

% \usepackage{needspace}


%%%%%%
 
\usepackage{listings}

\lstset{
  morekeywords={},
  sensitive=f,
  morecomment=[l]--,
  morestring=[d]",
  showstringspaces=false,
  basicstyle=\small\ttfamily,
  keywordstyle=\bf\small,
  commentstyle=\itshape,
  stringstyle=\sf,
  extendedchars=true,
  columns=[c]fixed
}

% CI-DESSOUS: conversion des caractères accentués UTF-8 
% en caractères TeX dans les listings...
\lstset{
  literate=%
  {À}{{\`A}}1 {Â}{{\^A}}1 {Ç}{{\c{C}}}1%
  {à}{{\`a}}1 {â}{{\^a}}1 {ç}{{\c{c}}}1%
  {É}{{\'E}}1 {È}{{\`E}}1 {Ê}{{\^E}}1 {Ë}{{\"E}}1% 
  {é}{{\'e}}1 {è}{{\`e}}1 {ê}{{\^e}}1 {ë}{{\"e}}1%
  {Ï}{{\"I}}1 {Î}{{\^I}}1 {Ô}{{\^O}}1%
  {ï}{{\"i}}1 {î}{{\^i}}1 {ô}{{\^o}}1%
  {Ù}{{\`U}}1 {Û}{{\^U}}1 {Ü}{{\"U}}1%
  {ù}{{\`u}}1 {û}{{\^u}}1 {ü}{{\"u}}1%
}

%%%%%%%%%%
% TAILLE DES PAGES (A4 serré)

\setlength{\parindent}{1cm}
\setlength{\parskip}{1ex}
\setlength{\textwidth}{17cm}
\setlength{\textheight}{22,7cm}
\setlength{\oddsidemargin}{-.7cm}
\setlength{\evensidemargin}{-.7cm}


\renewcommand{\labelenumi}{\arabic{enumi}.} 
\renewcommand{\labelenumii}{\arabic{enumi}.\arabic{enumii}}
\renewcommand{\labelenumiii}{\arabic{enumi}.\arabic{enumii}.\arabic{enumiii}}

%%%%%%%%%%

\newcommand\BackgroundPic{
\put(0,0){
\parbox[b][\paperheight]{\paperwidth}{%
\vfill
\includegraphics[width=\paperwidth,height=\paperheight,
keepaspectratio]{background.jpg}%
\vfill
}}}
%%%%%%%

\begin{document}

\AddToShipoutPicture{\BackgroundPic}


\renewcommand{\appendixtocname}{Annexes}
\DeclareGraphicsExtensions{.pdf,.png,.jpg}

\begin{titlepage}
\setlength{\parindent}{0cm}

\begin{center}

% Upper part of the page
 \begin{figure}[!h]
\includegraphics[bb=-550 -10 -250 20, scale=0.7]{./logo.pdf}
\end{figure}
% logo.pdf: 612x792 pixel, 72dpi, 21.59x27.94 cm, bb=0 0 612 792


\vspace{4cm}
\rule{\linewidth}{.5pt}
\vspace{2mm}


\begin{center}
{\LARGE GRAND CERCLE MOBILE - GCM}

\vspace{1cm}


{\Huge \bf DOCUMENT D'INSTALLATION}


\end{center}


\vspace{1cm}

%===================================================
\begin{center}
$ $\\
\large{ \textbf{Luiza CICONE - Jérémy KREIN - Jérémy LUQUET - Paul MAYER}}\\
\large{ \textbf{ISI - IF}}\\
\large{ \textbf{Professeur coordonnateur : Gaëlle CALVARY}}

$ $\\
\end{center}
\rule{\linewidth}{.5pt}


\vfill

% Bottom of the page

{\large Mai 2012}

\end{center}
\end{titlepage}

\tableofcontents

\newpage


\section{Installation de l'application iOS}
\subsection{Installation pour la maintenance}
Pour effectuer la maintenance de l'application iOS, il est nécessaire de posséder un Mac munit du logiciel Xcode, disponible sur l'Apple Store. 
Il est possible d'installer les différentes versions des simulateurs iOS pour pouvoir tester sur les différentes versions de l'OS (4.3, 5.0, 5.1). Il n'est plus possible de compiler l'application pour les versions plus anciennes d'OS.

Dans le dossier iOS se trouve le projet \lstinline!Grand Cercle.xcodeproj! qui s'ouvre automatiquement avec Xcode.

\subsection{Installation pour l'utilisateur}
L'application sera disponible sur l'Apple Store, donc les utilisateurs pourront la télécharger et l'installer comme tout autre application.

Une autre manière pour installer l'application est de connecter l'iPhone (3GS, 4 et 4S) ou l'iPod (4eme génération) à l'ordinateur avec un câble USB. Dans Xcode il faut choisir la Scheme le nom du dispositif et appuyer sur Run pour compiler et rouler le projet.
Dans ce cas, il est nécessaire d'avoir un profil de développeur Apple et d'associer dans un Provisioning Profile cet profil avec l'UDID du dispositif sur \texttt{http://developer.apple.com/devcenter/ios}. Après le télécharger, il faut l'ajouter en Xcode (se fait automatiquement en l'ouvrant). Dans le Target Grand Cercle, sous le Code Signing Identity il faut choisir le profil de développeur correspondant.


\section{Installation de l'application Android}
\subsection{Installation pour la maintenance}
Pour effectuer la maintenance de l'application Android, il est nécessaire de disposer du logiciel libre \texttt{eclipse}, disponible à l'adresse \href{http://www.eclipse.org/downloads/}{http://www.eclipse.org/downloads/} ainsi que le \texttt{plugin ADT} qui s'installe directement dans \texttt{eclipse} depuis le dépôt \href{http://dl-ssl.google.com/android/eclipse/}{http://dl-ssl.google.com/android/eclipse/}.

Pour tester directement l'application sur un ordinateur, le \texttt{plugin ADT} propose des émulateurs qui permettent de simuler des téléphones ou tablettes qui mettent en oeuvre différentes versions d'Android.

Cependant il est important de noter que les différents émulateurs proposés ne permettent pas de tester toutes les fonctionnalités de notre application. En effet, l'import d'un événement dans le calendrier du téléphone n'est pas géré par les émulateurs. Cette fonctionnalité doit donc être testée directement sur un téléphone sous Android. Ceci peut être réalisé en reliant directement le téléphone à l'ordinateur via un câble \texttt{USB}.

Pour pouvoir éditer le code source, il suffit d'importer le répertoire GrandCercleMobile dans un nouveau projet dans \texttt{eclipse}.
\subsection{Installation pour l'utilisateur}
L'utilisateur doit simplement installer le fichier \texttt{GrandCercleMobile.apk} pour disposer de l'application sur son téléphone.

Quand l'application sera disponible sur le \texttt{Play Store}, l'utilisateur pourra simplement la télécharger directement comme toute application Android.

\end{document}

